% ============================================================================
%  Section IV — Data-derived characterisation  (Measurement, Elsevier)
%  Passive voice / third person throughout. Requires: graphicx.
%  Figures: fig_psd_usable_band.png, fig_channel_mi.png (in results/).
%  Numbers from results/analysis_B.{md,json}.
% ============================================================================

\section{Data-derived characterisation of the signal channels}
\label{sec:charact}

Before any classifier is trained, it is instructive to establish what
fault-related information the recorded channels can, in principle, convey. It is
important to be precise about what follows and what does not. A metrological
characterisation of a MEMS accelerometer establishes sensitivity traceable to a
reference standard, amplitude linearity and frequency response, cross-axis
sensitivity, and thermal behaviour, and it requires controlled excitation and
reference instrumentation~\cite{prato2021mems, schiavi2023mems, xu2021mems,
wang2025mems, xu2021mems2, khankalantary2021mems}. No such experiments are
possible on an archived dataset, and none are claimed here.

What is presented instead is a \emph{data-derived characterisation of the signal
channels}: quantities computed directly from the delivered samples that bound
what those channels can carry. Three are reported---the observed minimum output
increment, an in-operation broadband power spectral density (PSD) floor, and the
band made observable by the sampling rate---together with the relative diagnostic
value of the raw and gravity-compensated triads. These are properties of the
delivered signal, not of the transducer: each conflates the sensor with the
device's internal processing, its acquisition software and the machine being
measured. They are nonetheless the quantities that bound the diagnostic task, and
they motivate the sensing choices analysed later.

\subsection{Observed output increment and in-operation signal floor}
The smallest positive difference between consecutive distinct sample values was
extracted from a healthy, unloaded recording. This gives approximately
$0.015$~mg for the raw acceleration channels and $0.001$~mg for the
gravity-compensated channels. The quantity is the granularity of the delivered
numerical stream, not the resolution of the transducer: the finer step of the
gravity-compensated triad reflects the additional arithmetic of the on-board
sensor fusion rather than a more sensitive sensing element, and neither figure
can be interpreted as a noise-limited resolution without the reference
measurements described above. An in-operation broadband floor was estimated from the
PSD of a healthy recording; the median PSD above
$30$~Hz was $8.2\times10^{-8}$~g$^2$/Hz ($-70.8$~dB re $1$~g$^2$/Hz). Because this
value is obtained on a \emph{running} machine, it includes residual mechanical
vibration and therefore upper-bounds, rather than isolates, the sensor's
intrinsic noise floor, which would require an acquisition with the device at
rest~\cite{srokosz2024mems, capuano2023smartphone}; no such recording exists in
this dataset. Even as an upper bound, however, it lies several orders of magnitude
below the fault-induced vibration, which reaches amplitudes on the order of
$0.1$--$0.2$~g (Fig.~\ref{fig:signals}); sensor resolution and noise are thus not
the limiting factor for the coarse fault-discrimination task considered here.

\subsection{Band made observable by the sampling rate}
At a sampling rate of $100$~Hz the observable spectrum is bounded by the Nyquist
frequency of $50$~Hz~\cite{nyquist1924certain}. This is a limit imposed by the
acquisition settings, not a measured property of the accelerometer, whose own
usable bandwidth is not accessible from these data.
Fig.~\ref{fig:psd} shows the Welch PSD of the gravity-compensated $g_{userx}$
channel for the six health states under identical operating conditions. Two
observations follow. First, discriminative structure is present within the
$0$--$50$~Hz band: the healthy machine exhibits a spectrum roughly an order of
magnitude lower in power than the fault conditions, and distinct spectral peaks
appear for the rotor and electrical faults. The dominant peaks near $24$~Hz
coincide with the shaft rotational frequency at $50$~Hz supply
(Table~\ref{tab:speeds}, e.g.\ $1425$--$1486$~r/min $\approx 23.8$--$24.8$~Hz),
confirming that the low-frequency band captures the rotational fundamental and
its low-order harmonics. Second, the characteristic bearing defect frequencies
(ball-pass and ball-spin frequencies and their harmonics) lie largely above
$50$~Hz at the higher supply frequencies---although some low-order components fall
within the band at $30$~Hz supply---and are, for the most part, not represented.
The dataset is therefore better suited to time-domain and low-frequency analysis
than to the high-frequency envelope analysis conventionally used for incipient
bearing diagnosis, and a corresponding difficulty in separating the fine
lubrication-severity states is anticipated.

One caution should be attached to this spectral reading. That the rotational
fundamental is present and visibly discriminative in the PSD does not establish
that a time-domain classifier makes use of it. Section~\ref{subsec:antialias}
shows the opposite: decimation that removes the fundamental entirely improves
leakage-free accuracy, so the descriptors used here are evidently not relying on
that component. The spectra below therefore delimit what the channel carries, not
what the classifier exploits.

\subsection{Raw versus gravity-compensated channels}
The dataset provides both the raw acceleration and the gravity-compensated
(\textit{userAcceleration}) triads, and their relative diagnostic value is not
self-evident. Two questions were separated: whether the on-device gravity
compensation removes vibration \emph{power}, and whether it removes diagnostic
\emph{information}. For the first, the alternating-current (AC) power of each raw
channel was compared with that of its gravity-compensated counterpart over every
recording; the mean raw-to-compensated power ratio was $1.00$ on all three axes
($0.999$, $1.001$, $1.000$ for $x$, $y$, and $z$). The compensation therefore
does not attenuate the vibration power. It does, however, \emph{reshape} it: the
mean-removed raw signal and the compensated signal are only moderately correlated
(Pearson $0.58$, $0.46$, and $0.46$ for $x$, $y$, and $z$), so the complementary
filter redistributes the low-frequency content---altering its phase and
correlation structure---rather than discarding it, and the two triads are
complementary rather than redundant.

For the second question, the mutual information between each channel's time-domain
feature set and the fault label was estimated (Fig.~\ref{fig:mi}). The raw
channels were consistently more informative (mean mutual information $0.69$) than
the gravity-compensated channels ($0.50$), a ratio of $1.39$. Because the raw
channels also retain a per-recording static offset that varies with condition, a
legitimate concern is that this advantage merely encodes recording identity---the
very leakage this work otherwise criticises. Two controls indicate that it does
not: removing the offset-dependent \emph{mean} feature leaves the ratio essentially
unchanged ($1.31$), and restricting the estimate to a DC-invariant feature subset
(standard deviation, variance, skewness, kurtosis, zero-crossing rate, and
amplitude entropy) yields a comparable ratio ($1.32$). The raw-channel advantage
is thus robust to the removal of static offsets and reflects genuine dynamic
content. The raw channels are therefore preferable for condition monitoring, and
this informativeness ordering motivates the channel-selection study reported in
Section~\ref{sec:results}.

% ---------------------------------------------------------------------------
\begin{figure}[t]
\centering
\includegraphics[width=\columnwidth]{fig_psd_usable_band.png}
\caption{Welch power spectral density of the $g_{userx}$ channel for the six
health states (loaded, $50$~Hz supply). The analysable band is bounded by the
$50$~Hz Nyquist frequency; the healthy state is markedly lower in power, and
rotor/electrical faults produce distinct low-frequency peaks.}
\label{fig:psd}
\end{figure}

% ---------------------------------------------------------------------------
\begin{figure}[t]
\centering
\includegraphics[width=\columnwidth]{fig_channel_mi.png}
\caption{Mean mutual information between each channel's time-domain feature set
and the fault label. The raw channels ($g_x$, $g_y$, $g_z$) are more informative
than the gravity-compensated channels ($g_{userx}$, $g_{usery}$, $g_{userz}$).}
\label{fig:mi}
\end{figure}
